\subsection{Exact thresholds for a supplied error certificate}

\begin{lemma}[Exact exponential blocking threshold]
    \label{lem:ldp_exact_exponential_blocking_threshold}
    \lean{MPSPreparation.mul_pow_mul_exp_neg_le_iff,
        MPSPreparation.mul_pow_mul_exp_neg_nat_le_iff,
        MPSPreparation.mul_div_mul_exp_neg_le_iff}
    \notready
    Let $K,N,\varepsilon,r>0$ and $k\ge0$ be an integer. Then
    \[
        KN^k e^{-rq}\le\varepsilon
        \quad\Longleftrightarrow\quad
        q\ge\frac{\log K+k\log N-\log\varepsilon}{r}.
    \]
    The least nonnegative integer satisfying this certificate is the
    nonnegative ceiling of the displayed threshold. If the certificate
    retains the number-of-blocks factor and $q>0$, then
    \[
        K\frac Nq e^{-rq}\le\varepsilon
        \quad\Longleftrightarrow\quad
        rq+\log q\ge\log K+\log N-\log\varepsilon.
    \]
\end{lemma}
\begin{proof}
    Express the left sides as exponentials, use the strict monotonicity
    of the exponential and rearrange. For integral scales, apply the
    defining property of the nonnegative ceiling.
\end{proof}

These equivalences are exact for the stated numerical certificates.
They do not establish such a certificate for an inhomogeneous sequence,
nor an optimal exponent for its actual approximation errors.
The sampled decay in (S42) of~\cite{Malz2024Preparation} does not supply
a uniform sharp rate for the unspecified long-range sequence class.

\subsection{An exact tree for the actual inhomogeneous chain}

\begin{lemma}[Concatenation of ordered chain products]
    \label{lem:ldp_chain_eval_append}
    \lean{MPSChainTensor.eval_append}
    \notready
    For two site-dependent chains $A,B$ and physical words $\sigma,\tau$,
    the ordered matrix product along their concatenation satisfies
    $F_{AB}(\sigma\tau)=F_A(\sigma)F_B(\tau)$.
\end{lemma}
\begin{proof}
    Concatenate the ordered lists of local matrices and multiply.
\end{proof}

\begin{definition}[Contraction of child remainders]
    \label{def:ldp_chain_pair_merge}
    \lean{MPSPreparation.chainPairMerge}
    \notready
    Identify the virtual space with pairs $(a,b)$ of indices in
    $\{1,\ldots,D\}$. For $D^2\times D^2$ matrices $P,Q$, set
    \[
        C_{(u,v),(a,b)}
        =\sum_{c=1}^D P_{u,(a,c)}Q_{v,(c,b)}.
    \]
\end{definition}

\begin{lemma}[The actual parent lies in the child ranges]
    \label{lem:ldp_chain_pair_append_factorization}
    \lean{MPSPreparation.chainPairMatrix_append_eq_joinMatrix_mul}
    \notready
    \uses{lem:ldp_chain_eval_append, def:ldp_chain_pair_merge}
    Let $F_A,F_B$ be the physical matrices of two site-dependent chains.
    If $F_A=XP$ and $F_B=YQ$, then the physical matrix of their
    concatenation satisfies $F_{AB}=(X\otimes Y)C$, where $C$ is the
    contraction of $P,Q$ from Definition~\ref{def:ldp_chain_pair_merge}.
    The rows of the tensor product are read in physical chain order.
\end{lemma}
\begin{proof}
    At the column $(a,b)$ and word $\sigma\tau$, expand
    $(F_A(\sigma)F_B(\tau))_{ab}$ over the common bond index.
    Substitute the two child factorizations and exchange the finite sums.
\end{proof}

\begin{lemma}[An isometric merge for actual adjacent intervals]
    \label{lem:ldp_chain_pair_actual_merge}
    \lean{MPSPreparation.exists_isometry_chainPairMatrix_append}
    \notready
    \uses{lem:ldp_chain_pair_append_factorization,
        lem:ldp_isometric_polar_merge}
    Let $D\ge1$. If $X,Y$ are isometries with
    $XP_A=F_A$ and $YP_B=F_B$, where $P_A,P_B$ are the positive polar
    factors of the child physical matrices, there is an isometry
    $W\colon\C^{D^2}\to\C^{D^4}$ satisfying
    $(X\otimes Y)WP_{AB}=F_{AB}$.
\end{lemma}
\begin{proof}
    The preceding lemma establishes that $F_{AB}$ lies in the range of
    $X\otimes Y$. Complete its polar factor within that range by
    Lemma~\ref{lem:ldp_isometric_polar_merge}.
\end{proof}

\begin{lemma}[Agreement on the polar support]
    \label{lem:ldp_factor_agrees_on_polar_support}
    \lean{Matrix.mul_polarSupport_eq_polarIso_of_mul_polarPos_eq}
    \notready
    If $M=VP$ is the canonical polar decomposition with support $\Pi$
    and $XP=M$, then $X\Pi=V$. No injectivity or isometry of $X$ is needed.
\end{lemma}
\begin{proof}
    Choose $R$ with $PR=\Pi$. Then
    $X\Pi=XPR=MR=VPR=V\Pi=V$.
\end{proof}

\begin{lemma}[Exact polar tree on a feasible binary geometry]
    \label{lem:ldp_inhomogeneous_exact_polar_tree}
    \lean{MPSPreparation.exists_isometryTree_mul_polarPos_eq}
    \notready
    \uses{lem:ldp_chain_pair_actual_merge,
        lem:ldp_factor_agrees_on_polar_support,
        lem:ldp_isometric_polar_completion}
    Let $D\ge1$ and suppose a binary isometric tree of height $h$,
    virtual dimension $D^2$ and leaves of at most $R$ physical sites exists
    on a chain of length $n$. For every site-dependent chain on those
    $n$ sites, there is such a tree with physical map $X$ satisfying
    $XP=F$ and $X\Pi=V$ for the physical matrix $F$ and its canonical
    polar factors. The maps may depend on their intervals.
\end{lemma}
\begin{proof}
    Induct over the supplied binary geometry. At a leaf, its feasible
    dimension permits an isometric polar completion of the actual leaf
    matrix. At a parent, split the actual chain into its two intervals,
    apply the induction hypotheses and merge by
    Lemma~\ref{lem:ldp_chain_pair_actual_merge}. Agreement on the polar
    support follows from Lemma~\ref{lem:ldp_factor_agrees_on_polar_support}.
\end{proof}

\begin{lemma}[Balanced geometry and exact inhomogeneous polar tree]
    \label{lem:ldp_balanced_inhomogeneous_polar_tree}
    \lean{MPSPreparation.IsometryTree.nonempty_of_leaf_dimension,
        MPSPreparation.exists_isometryTree_chainPairMatrix}
    \notready
    \uses{lem:ldp_inhomogeneous_exact_polar_tree}
    For $s\ge1$ and $d^s\ge D^2\ge1$, every site-dependent chain of
    length $s2^h$ has an exact polar tree of height $h$, virtual dimension
    $D^2$, and leaves of at most $s$ sites.
\end{lemma}
\begin{proof}
    Choose an isometric inclusion into each $s$-site leaf and an
    isometry from $\C^{D^2}$ to $\C^{D^4}$ at every parent. This supplies
    a balanced feasible geometry. Apply
    Lemma~\ref{lem:ldp_inhomogeneous_exact_polar_tree}.
\end{proof}

\begin{lemma}[Coherent measurement rounds for position-dependent tree gates]
    \label{lem:ldp_interval_tree_measurement_rounds}
    \lean{MPSPreparation.exists_rounds_blockLayerOp_treeUnitaryOp}
    \notready
    Fix a physical dimension, a register width and a bound on the physical
    leaf lengths. Arbitrary unitary node and leaf gates on disjoint binary
    layouts have a measurement implementation of depth at most $C(h+1)$,
    where $C$ is independent of the number of blocks, their lengths and
    their gates. Each gate may depend on its block and vertex.
    On inputs with zero scratch registers, every outcome history acts as
    one scalar times the full tree gate product; that scalar is independent
    of the input and hence shared by all coherent sectors.
\end{lemma}
\begin{proof}
    Implement the register levels by the coherent teleportation rounds
    and the bounded leaves by one local circuit. Compose the rounds.
    Their unitary target product makes each outcome history a scalar
    multiple of the same linear map on the zero-scratch subspace.
\end{proof}

The exact factorization and the gate implementation remain distinct.
To combine them, choose the extensions of each interval isometry and
identify the entire placed gate product on encoded inputs with the
algebraic tree map, including every scratch site. An approximation rate
for a specified sequence class is a separate requirement.
